% Listening-examples table: BASIC mode. Hyperlinks to MP3s hosted on
% GitHub Pages (paper-audio/ in the repo). Companion to
% listening_examples_advanced_table.tex - same 4 settings, same
% unprocessed reference, Basic (full-band) ducking instead of Advanced.
% Requires packages already in main.tex's preamble: booktabs, siunitx,
% hyperref (colorlinks already configured there).
%
% All 5 files have +3dB gain applied UNIFORMLY (including unprocessed) so
% differences a listener hears reflect the actual processing, not a level
% difference between clips - see paper-audio/'s commit messages for the
% clipping-headroom checks (no file exceeds -1.41dBFS after the boost).
%
% Row labels (Weakest/Best/2nd best) intentionally match Table~\ref{tab:param-sweep}'s
% rank numbers, which were derived from ADVANCED-mode margin gain - kept
% consistent across both listening tables so a reader can cross-reference
% one settings combination across all three tables. Confirmed these 4
% settings rank in the same relative order under Basic's own margin gain
% too (Run 1 > Run 2 > Run 3 > Run 9 for both modes), so the labels aren't
% misleading applied here.
%
% Base URL: https://jmicensky.github.io/sidechain-keyed-unmasking-demo/paper-audio/

\noindent Table~\ref{tab:listening-examples-basic} links to audio examples
for Basic (full-band) ducking: the unprocessed reference, and Basic mode
under the same four representative settings used in
Table~\ref{tab:listening-examples-advanced}'s Advanced-mode examples,
letting the two be compared directly setting-for-setting.

\begin{table}[h]
\centering
\footnotesize
\caption{Listening examples, Basic (full-band) ducking: the same
$\sim$\SI{82}{\second} Construction Scene / Dialogue-priority clip,
unprocessed and under four settings from the parameter sweep. All five
files have \SI{3}{\decibel} of gain applied uniformly for comfortable
playback level; audible differences between clips reflect the processing,
not a level mismatch.}
\label{tab:listening-examples-basic}
\begin{tabular}{@{}l S[table-format=-2.0] S[table-format=1.1] S[table-format=2.0] S[table-format=2.0] S[table-format=3.0] l@{}}
\toprule
\textbf{Condition} & {\textbf{Thr.} (\si{\decibel})} & {\textbf{Ratio}} & {\textbf{Knee} (\si{\decibel})} & {\textbf{Atk.} (\si{\milli\second})} & {\textbf{Rel.} (\si{\milli\second})} & \textbf{Listen} \\
\midrule
Unprocessed & \multicolumn{5}{c}{(dry reference, no ducking)} & \href{https://jmicensky.github.io/sidechain-keyed-unmasking-demo/paper-audio/unprocessed.mp3}{Listen} \\
Weakest (Table~\ref{tab:param-sweep} \#9) & -24 & 4.0 & 6 & 5 & 400 & \href{https://jmicensky.github.io/sidechain-keyed-unmasking-demo/paper-audio/basic_th-24_r4_k6_a5_rel400.mp3}{Listen} \\
Best (Table~\ref{tab:param-sweep} \#1) & -36 & 4.0 & 20 & 5 & 400 & \href{https://jmicensky.github.io/sidechain-keyed-unmasking-demo/paper-audio/basic_th-36_r4_k20_a5_rel400.mp3}{Listen} \\
2nd best (Table~\ref{tab:param-sweep} \#2) & -48 & 1.5 & 12 & 5 & 400 & \href{https://jmicensky.github.io/sidechain-keyed-unmasking-demo/paper-audio/basic_th-48_r1.5_k12_a5_rel400.mp3}{Listen} \\
\#2, faster release & -48 & 1.5 & 12 & 5 & 120 & \href{https://jmicensky.github.io/sidechain-keyed-unmasking-demo/paper-audio/basic_th-48_r1.5_k12_a5_rel120.mp3}{Listen} \\
\bottomrule
\end{tabular}
\end{table}
